% TOP_PROBLEM: {"id": 3700033, "title": "Carleson–Jones quarter conjecture", "category_id": 8, "category": {"id": 8, "name": "analysis", "display_name": "Analysis", "description": "Limits, continuity, calculus, and function theory.", "slug": "analysis", "order_index": 8, "created_at": "2026-07-31T15:26:25.671Z"}, "status": "open", "problem_number": "AMR-036-0033", "set_id": 13}
% FIRST_POSTED: 2026-10-01
% ENTRY_KIND: statement_only
% UnsolvedMath ID 3700033; source catalogue code AMR-036-0033
% !TeX program = lualatex
% Definition and sources only; this file is not an LLM solution attempt.
\documentclass[11pt,a4paper]{article}
\usepackage[margin=25mm]{geometry}
\usepackage{fontspec}
\setmainfont{lmroman10-regular.otf}[BoldFont=lmroman10-bold.otf,
 ItalicFont=lmroman10-italic.otf,BoldItalicFont=lmroman10-bolditalic.otf]
\usepackage{amsmath,amssymb,mathtools}
\usepackage{unicode-math}
\setmathfont{latinmodern-math.otf}
\AtBeginDocument{\let\setminus\smallsetminus}
\usepackage{xurl,hyperref}
\hypersetup{colorlinks=true,urlcolor=blue}
\setcounter{secnumdepth}{0}
\setlength{\parindent}{0pt}
\setlength{\parskip}{5pt}
\setlength{\emergencystretch}{3em}
\begin{document}
\section*{Carleson–Jones quarter conjecture}\label{3700033}
{\small UnsolvedMath ID 3700033 · AMR-036-0033 · Analysis · AMR Open Problem Lists}
\subsection{Definitions and mathematical statement}
Let $E\subset\mathbb C$ be a connected compact set
of positive logarithmic capacity. Let $\Omega_E$
be the unbounded component of its complement.
Its Green function with pole at infinity is
the positive harmonic function $G_E$ on
$\Omega_E$ normalized by
\[
 G_E(z)=\log|z|-\log\operatorname{cap}(E)+o(1)
                                      \quad(z\to\infty).
\]
Require regularity: $G_E$ extends continuously by
zero to the boundary. Extend it by zero also
over $E$ and the bounded complementary components.

For $\varepsilon>0$, define the total level-set
length and its growth exponent by
\[
 \ell_E(\varepsilon)=
             \mathcal H^1\{z:G_E(z)=\varepsilon\},
 \qquad
 \beta_E=\limsup_{\varepsilon\downarrow0}
             \frac{\log\ell_E(\varepsilon)}{-\log\varepsilon}.
\]
Length means one-dimensional Hausdorff measure,
including all components of the level curve.
Put $\beta=\sup_E\beta_E$, over all such
regular connected compact sets.

Is $\beta=1/4$? A proof requires a universal
upper bound $\beta_E\leq1/4$ and examples
whose exponents approach $1/4$ arbitrarily
closely. It need not supply a single compact
set attaining the supremum.
The limiting variable is the Green-function
level, rather than Euclidean distance to the
compact set. Translation or fixed rescaling
of $E$ changes level lengths only by a
constant factor and leaves this exponent
unchanged. No smoothness or finite boundary
length is assumed for $E$ itself.
\subsection{Short English statement}
Is one quarter the largest possible growth exponent for lengths of Green-function level curves around a connected compact set?
\subsection{Sources}
\begin{itemize}
\item[S1] ULAM.AI and the UnsolvedMath Contributors, supplied problems.json, record 3700033 (AMR-036-0033), AMR Open Problem Lists. Catalogue identity and source transcription; CC BY 4.0.
\url{https://huggingface.co/datasets/ulamai/UnsolvedMath}.
\item[S2] Eremenko - My favorite unsolved problems, Littlewood constants, Section 2 conjecture. Original source indexed by AMR-036-0033.
\url{https://www.math.purdue.edu/~eremenko/uns1.html}.
\item[S3] A. Eremenko, Some constants coming from the work of Littlewood, Section 2. Author-maintained primary problem note.
\url{https://www.math.purdue.edu/~eremenko/dvi/lit.pdf}.
\end{itemize}
\end{document}
